%% Full proofs for terms, identity consequence, and free algebras.

\section{Terms, equational completeness, and free algebras}

\begin{annexitem}{annex:p10.eq.1}{thm:p10-term-universal}{Universal property and substitution law for terms}
\begin{proof}
The recursive evaluation clauses define a function on every finite
term. They give \(v(x)\) on generators and preserve each operation
by their defining equation, including nullary operations.
Thus the extension is a homomorphism.
Any homomorphism extending \(v\) has those values on variables and
constants. Induction on a compound term forces its value from the
values of its immediate subterms, proving uniqueness.

For the substitution law, a variable \(x\) gives
\((x\sigma)^A(v)=(\sigma(x))^A(v)=w(x)\).
For a compound term, apply the induction hypothesis to each
subterm and then apply the common operation \(f^A\).
For a nullary term both sides equal the designated constant of \(A\).
This proves the law for all terms.

Finally, homomorphisms commute with term evaluation.
For variables this is the definition of \(h\circ v\); for a basic
compound term it follows from the induction hypotheses and
preservation of the basic operation by \(h\).
The nullary case uses preservation of constants.
\end{proof}
\end{annexitem}

\begin{annexitem}{annex:p10.eq.2}{thm:p10-equational-completeness}{Completeness of equational logic}
\begin{proof}
\paragraph{Soundness.}
Fix a nonempty algebra \(A\models\Gamma\).
Induct on a finite derivation. Premises hold as identities by
assumption. Reflexivity, symmetry, and transitivity preserve equality,
and compatibility preserves it under every basic operation.
For substitution, suppose \(A\models s\approx t\), and fix
any valuation \(v:V\to A\). Define
\[
 w(x)=(\sigma(x))^A(v).
\]
The substitution law gives
\[
 (s\sigma)^A(v)=s^A(w)=t^A(w)=(t\sigma)^A(v).
\]
The middle equality uses the identity at the valuation \(w\).
Since \(v\) was arbitrary, the substituted equation is an identity
of \(A\). This proves soundness.

\paragraph{The deductive congruence.}
First compare operation compatibility with contextual replacement.
From a derived \(s\approx t\), replacement at the whole term is
that same derivation.
If the selected occurrence lies in one argument of
\(f(u_1,\ldots,u_n)\), apply the induction hypothesis inside that
argument, use reflexivity on every other argument, and apply
compatibility. Induction on the finite distance from the selected
occurrence to the root proves replacement in any term context.
Conversely, given all \(s_i\approx t_i\), replace the arguments of
\(f(s_1,\ldots,s_n)\) one at a time and compose the resulting
equations by transitivity. This gives operation compatibility.
For arity zero the identity is reflexivity.
Thus these two presentations have the same deductions.

Reflexivity, symmetry, and transitivity make \(\theta_\Gamma\)
an equivalence relation; compatibility makes it a congruence.
Every endomorphism of the term algebra is induced by its values on
the variables, by Theorem~\ref{thm:p10-term-universal}.
The substitution rule therefore makes \(\theta_\Gamma\) fully
invariant.
It contains every premise pair.
Conversely, any fully invariant congruence containing those pairs
contains every derivable pair, by induction on the derivation:
its equivalence laws, compatibility, and full invariance verify
exactly the corresponding inference rules.
Thus \(\theta_\Gamma\) has the claimed leastness property.

Let \(B=T_\tau(V)/\theta_\Gamma\).
Congruence compatibility makes the operations
\[
 f^B([t_1],\ldots,[t_n])=[f(t_1,\ldots,t_n)]
\]
well defined. Since the term algebra is nonempty, so is \(B\).

\paragraph{Every premise holds under every quotient valuation.}
Fix \(p\approx q\in\Gamma\) and a valuation \(v:V\to B\).
Only finitely many variables occur in \(p,q\).
For each of them choose a representative term \(r_x\) with
\([r_x]=v(x)\).
Fix any term \(r_0\in T_\tau(V)\), and extend these finitely many
choices to a substitution by setting \(\sigma(x)=r_0\) on all
other variables.
Induction on the subterms of \(p,q\) gives
\[
 p^B(v)=[p\sigma],\qquad q^B(v)=[q\sigma].
\]
The premise and substitution rules give
\(\Gamma\vdash_{\mathrm{eq}}p\sigma\approx q\sigma\), so these
classes are equal.
Hence \(B\models p\approx q\); since the premise was arbitrary,
\(B\models\Gamma\).
Only a finite family of representatives is selected for each
equation. This argument does not require a simultaneous choice of
representatives for the whole quotient.

\paragraph{The countermodel and completeness.}
For the canonical valuation \(\eta(x)=[x]\), induction gives
\(u^B(\eta)=[u]\) for every term \(u\).
If \(\Gamma\nvdash_{\mathrm{eq}}s\approx t\), then
\((s,t)\notin\theta_\Gamma\), so \(s^B(\eta)\ne t^B(\eta)\).
The algebra \(B\) satisfies all premises as identities and fails
the conclusion. This proves the contrapositive of completeness.

Finally, a finite derivation uses only finitely many equations
from \(\Gamma\).
Collect them into \(\Gamma_0\).
The same derivation is a derivation from \(\Gamma_0\), and soundness
gives \(\Gamma_0\models_{\mathrm{alg}}s\approx t\).
Combining this observation with completeness proves finitariness,
even for an infinite premise set.
\end{proof}
\end{annexitem}

\begin{annexitem}{annex:p10.eq.3}{thm:p10-free-algebra}{Construction and uniqueness of free algebras}
\begin{proof}
For each algebra \(A\in\mathcal V\) and map \(v:X\to A\), the kernel
of term evaluation is a congruence.
The relation \(\theta_{\mathcal V,X}\), defined by requiring
membership in every such kernel, is therefore a congruence:
the equivalence laws and compatibility hold in each kernel and
hence hold under their common condition.
This defines a relation on the set \(T_\tau(X)^2\), even when the
class of models \(\mathcal V\) is proper.
Let \(F=T_\tau(X)/\theta_{\mathcal V,X}\).

\paragraph{The quotient belongs to the variety.}
Fix \(p\approx q\in\Gamma\), with variables in the declared identity
alphabet \(V\), and a valuation \(v:V\to F\).
Choose representative terms \(r_x\in T_\tau(X)\) for the values
of the finitely many variables in \(p,q\).
Using a fixed term elsewhere gives a substitution
\(\sigma:V\to T_\tau(X)\).
As in the preceding proof,
\[
 p^F(v)=[p\sigma],\qquad q^F(v)=[q\sigma].
\]
For every \(A\in\mathcal V\) and \(u:X\to A\), the substitution law
and \(A\models p\approx q\) give
\[
 (p\sigma)^A(u)=p^A(x\mapsto r_x^A(u))
              =q^A(x\mapsto r_x^A(u))=(q\sigma)^A(u).
\]
Here \(r_x=\sigma(x)\) for variables outside the finite list as well.
Thus \((p\sigma,q\sigma)\in\theta_{\mathcal V,X}\).
The displayed quotient values are equal.
This holds for every premise and valuation, so \(F\models\Gamma\)
and \(F\in\mathcal V\).

\paragraph{Existence and uniqueness of extensions.}
Given \(A\in\mathcal V\) and \(v:X\to A\), define
\[
 \bar v([t])=t^A(v).
\]
The definition of \(\theta_{\mathcal V,X}\) makes this independent
of the representative \(t\).
Evaluation preserves operations, so \(\bar v\) is a homomorphism,
and \(\bar v(\eta_X(x))=v(x)\).
Every element of \(F\) is the class of a term.
Its value under a homomorphism is determined by the generator
values and the interpretations of nullary operations.
Induction on terms therefore proves uniqueness, also when
\(X=\varnothing\) and all terms are ground.

If \((F,\eta)\) and \((F',\eta')\) are both free on \(X\), their
universal properties give homomorphisms
\(h:F\to F'\), \(k:F'\to F\) with
\(h\eta=\eta'\) and \(k\eta'=\eta\).
Both \(kh\) and the identity of \(F\) extend \(\eta\), so uniqueness
gives \(kh=1_F\).
Similarly \(hk=1_{F'}\).
These maps are inverse isomorphisms, and the extension uniqueness
also makes the isomorphism commuting with the generators unique.
\end{proof}
\end{annexitem}
